\documentclass[11pt]{article}
\usepackage[utf8]{inputenc}
\usepackage{geometry}
\geometry{a4paper, margin=1in}
\usepackage{titlesec}
\usepackage{enumitem}
\usepackage{xcolor}

\title{\textbf{Strategic Diagnostic Snapshot™}}
\author{RankPilot Intelligence}
\date{\today}

\begin{document}

\maketitle

\section*{I. Executive Summary}
\textbf{Firm:} González de Araujo Consultores \\
\textbf{Practice Area:} Data Protection \\
\textbf{Detected Practice Model:} Standard

\section*{II. Structural Positioning}
\textbf{Current Tier Assessment:} Unranked / New Entry \\
\textbf{Strategic Confidence:} 100.0\% \\
\textbf{Advantage:} Standard Market Offerings

\section*{III. Portfolio Analysis (Key Matters)}
\begin{itemize}
\item \textbf{Grupo Hermes Data Protection Framework} (Client: Grupo Hermes | Value: N/A)\\ The Grupo Hermes, Mexican diversified infrastructure and services conglomerate, Data Protection Framework demonstrates the firm’s central role in converting data-protection governance into documented, operational procedures with continuing oversight. The mandate was significant not merely because it concerned privacy documentation, but because it required the design and implementation of a thorough data protection and governance framework capable of addressing the client’s personal-data handling through connected governance measures.

For Grupo Hermes, developing privacy notices, implementing policies, and establishing ARCO rights management procedures were integral components of the framework designed and implemented by the firm. The team mapped personal data flows as the factual foundation for identifying how personal data moved through the client’s operations and for aligning the ensuing governance measures with those flows. Building on that mapping exercise, the firm developed privacy notices, implemented policies, and established ARCO rights management procedures, bringing together notice requirements, internal policies, and processes for managing ARCO rights within the same data protection and governance framework.

Arturo González de Araujo and Ernesto Zavala led the engagement. Their role encompassed the design and implementation of the thorough data protection and governance framework, as well as the central work of mapping personal data flows, developing privacy notices, implementing policies, and establishing ARCO rights management procedures. The firm’s subsequent appointment as ongoing data protection auditor and strategic advisor confirms a continuing role beyond the initial framework implementation, focused on auditing the operation of the data protection arrangements and providing strategic advisor support in relation to the framework.
\item \textbf{Mega Direct Data Protection Advisory} (Client: MEGA DIRECT, S.A. de C.V. | Value: N/A)\\ Representing MEGA DIRECT, S.A. de C.V., a customer experience, call center, and marketing services provider, in legal proceedings concerning the full lifecycle of data acquisition, processing, storage, and protection, this mandate demonstrates the submission’s evidence of contentious and operationally significant personal-data work. The engagement was not confined to abstract data-protection advice: it placed the client’s database-related activities at the centre of legal proceedings and required representation directed at preserving the continuity of those activities.

MEGA DIRECT, S.A. de C.V. instructed the team on complex data protection and database-related matters arising in a highly regulated environment. The work addressed the legally sensitive connection between the acquisition of data, its subsequent processing and storage, and its protection, with each element forming part of the proceedings in which the client was represented. By addressing those interrelated issues through legal representation, the team ensured business continuity for MEGA DIRECT, S.A. de C.V. while the relevant data and database-related matters were under consideration.

Arturo González de Araujo, Ernesto Zavala, and Damián Ortega led the mandate. Their involvement evidences a senior-led response to proceedings in which the client’s continued ability to manage data acquisition, processing, storage, and protection was directly relevant to business continuity in a highly regulated environment.
\item \textbf{Biocodex Data Protection Framework} (Client: Biocodex | Value: N/A)\\ Establishing a Personal Data Protection Department gave this mandate significance beyond a discrete compliance exercise: it placed personal data protection within an identifiable internal function while addressing the protection of sensitive data under Mexican regulations.

Biocodex, a global pharmaceutical company focused on probiotics and specialty healthcare, instructed the team in connection with the design and implementation of a structured personal data protection framework. The engagement required the team to translate the applicable Mexican regulatory requirements into an operational framework capable of supporting the client’s treatment of personal data and the particular protection required for sensitive data. Rather than being confined to a high-level assessment, the work was directed to the design and implementation of the structured framework itself, linking the regulatory objective of compliance with the concrete institutional step of establishing a Personal Data Protection Department.

Arturo González de Araujo, Ernesto Zavala and Damián Ortega led the matter. Their involvement connected the framework-design work with the establishment of the Personal Data Protection Department, ensuring that the department’s creation formed part of the same personal-data-protection exercise. The resulting structure supported Biocodex’s compliance with Mexican regulations and created a dedicated departmental basis for the protection of sensitive data.
\item \textbf{Hotel Riazor Data Protection Enhancement} (Client: Hotel Riazor | Value: N/A)\\ The significance of Hotel Riazor, business hotel offering lodging and event services in Mexico City’s Data Protection Enhancement mandate lies in its treatment of personal-data governance as an operational and compliance priority, linking the strengthening of the client’s internal data-protection arrangements to its broader operational efficiency and compliance programme.

Hotel Riazor instructed the team to advise on strengthening its personal data protection and governance framework. The mandate centred on reinforcing the Personal Data Department and on translating the governance objective into implemented policies and procedures for lawful data handling. The work was therefore directed not only at the internal departmental structure responsible for personal data, but also at embedding lawful data-handling procedures within the client’s broader operational efficiency and compliance programme.

The concrete result recorded for the engagement was the reinforcement of the Personal Data Department and the implementation of policies and procedures for lawful data handling, integrated into the broader operational efficiency and compliance programme. As stated in the underlying record: Hotel Riazor implemented policies and procedures for lawful data handling, integrated into the broader operational efficiency and compliance programme.

Arturo González de Araujo, Ernesto Zavala and Damián Ortega were the Lead Partners on the matter. Their involvement in advising on the strengthened personal data protection and governance framework, reinforcing the Personal Data Department, and implementing policies and procedures for lawful data handling demonstrates partner-led work connecting data protection with Hotel Riazor’s operational efficiency and compliance programme.
\item \textbf{Grupo Excelsior Data Protection Regularisation} (Client: Grupo Excelsior | Value: N/A)\\ Regularisation and restructuring of operations can expose whether data protection is treated as an isolated compliance requirement or incorporated into the organisation’s governing arrangements. In this mandate, the work for Grupo Excelsior, one of Mexico’s leading dairy producers, centred on bringing data protection into the regularisation and restructuring of operations, giving the exercise a clear operational and corporate-governance dimension rather than limiting it to a standalone policy review.

Grupo Excelsior engaged the team to advise on the regularisation and restructuring of operations with a focus on data protection. The firm’s role was to translate that focus into implemented internal measures: policies, protocols and training. Those measures addressed the practical task of embedding a data protection culture, aligning operations with regulations, and integrating data protection into corporate governance. The implementation outcome is reflected in the following formulation: Grupo Excelsior, protocols, and training to embed data protection culture, align with regulations, and integrate into corporate governance.

Arturo González de Araujo, Ernesto Zavala and Damián Ortega led the work. Their contribution connected the regularisation and restructuring of operations to the implementation of policies, protocols and training, ensuring that data protection was embedded as part of Grupo Excelsior’s corporate-governance arrangements. The resulting measures embedded a data protection culture, aligned the relevant operations with regulations, and integrated data protection into corporate governance.
\item \textbf{Grupo Modelquipo Data Protection Framework} (Client: Grupo Modelquipo | Value: N/A)\\ Grupo Modelquipo’s, an industrial equipment and machinery solutions provider, data protection framework mandate is significant because it placed the identification and mitigation of data-processing risks within the client’s corporate governance strategy, rather than treating privacy compliance as a standalone documentation exercise. The engagement therefore provides concrete evidence of work connecting data-processing risk management, internal governance and day-to-day compliance implementation.

Arturo González de Araujo, Ernesto Zavala and Damián Ortega advised on the data protection framework as part of Grupo Modelquipo’s corporate governance strategy. Their specific role encompassed identifying data-processing risks and mitigating those risks through practical compliance measures. The work moved from risk identification to implementation, addressing both the privacy information to be provided and the internal procedures required to support compliance.

The mandate resulted in the implementation of privacy notices and protocols, together with training for compliance. In the client-description wording supplied for this matter, the outcome must be recorded as follows: “Grupo Modelquipo implemented privacy notices and protocols, delivered training for compliance”. This combination of risk identification and mitigation, privacy notices, protocols and compliance training established the data protection framework as an operational component of Grupo Modelquipo’s corporate governance strategy.
\item \textbf{Tiendas Chedraui Data Protection Advisory} (Client: Tiendas Chedraui | Value: N/A)\\ Tiendas Chedraui, a Mexican supermarket and retail chain, Data Protection Advisory records the application of data-protection analysis to the client’s planned new store openings and expansion, placing the legal assessment at the point where new operational activity gives rise to additional data-processing considerations. The work is recorded under the client descriptor: Tiendas Chedraui, reviewed documentation and compliance with data protection obligations.

Arturo González de Araujo and Damián Ortega led the mandate for Tiendas Chedraui. Their specific role was to conduct a legal assessment of data protection in connection with the new store openings and expansion, rather than limiting the engagement to a general review of internal materials. The team evaluated risks associated with data processing arising from that planned activity. In parallel, Arturo González de Araujo and Damián Ortega reviewed the documentation connected with those operations and assessed compliance with data protection obligations.

The significance of the mandate lies in the way it connects three defined elements of the client’s expansion activity: the legal assessment of data protection for new store openings and expansion; the evaluation of risks associated with data processing; and the review of documentation and compliance with data protection obligations. The resulting work records a defined data-protection review directed to the documentation and processing risks associated with Tiendas Chedraui’s new store openings and expansion.

\end{itemize}

\section*{IV. Strategic Evolution Path}
\begin{itemize}
\item Maintain current strategy.
\end{itemize}

\vfill
\centering
\small{This document is a structural assessment and does not guarantee ranking results.}

\end{document}